\scanpage{9}
\mspage{4}
quels nous allons nous borner par la suite -- on a donc des homs canoniques
\[
(6)\qquad B_K\longrightarrow B_{\mathbb Q}\qquad\Pi_K\longrightarrow\Pi_{\mathbb Q}
\]
qui s'explicitent, quand on a choisi un objet de $\Pi_K$ i.e. une $\bar K/K$, d'où un $\overline{\mathbb Q}/\mathbb Q$, par un hom de groupes profinis
\[
(7)\qquad E_{\bar K/K}\longrightarrow\Gamma_{\overline{\mathbb Q}/\mathbb Q}.
\]
Par la suite, on regardera toujours $B_K$, $\Pi_K$ ou $E_{\bar K/K}$ comme munis de cette structure supplémentaire -- \unsure{ce sont} les morphismes (de topos, de groupoïdes, ou de groupes profinis) ``\uline{arithmétiques}'', \unsure{domineront|dominent} la situation. !

Un intérêt particulier s'attache au noyau de (7), que je note $\pi_{\bar K/K}$ -- on l'appelle ``partie géométrique'' du groupe Galois $E_{\bar K/K}$, par opposition au quotient
\[
E_{\bar K/K}/\pi_{\bar K/K}=\Gamma^{\deleted{\unsure{a}}}_{\bar K/K}\hookrightarrow\Gamma_{\overline{\mathbb Q}/\mathbb Q},
\]
\margin{on va noter $\Gamma=\Gamma^{\deleted{\unsure{a}}}_{\bar K/K}$ cette ``partie arithmétique''}
que j'appelle sa
\margin{8}
